\subsection{Magnetization concentration and bounded
covariance}\label{proof:thin_shell_obstruction:magnetization-concentration-and-bounded-covariance}

The spin means vanish. Under \(\nu _n\), \(\mathbb{E} S^{2}=16\), and
exchangeability gives

\[
\mathbb E_{\nu_n}X_iX_j=\frac{16-n}{n(n-1)}\quad(i\ne j).
\]

The covariance eigenvalues are \(16/n\) in the all-ones direction and
\((n^{2}-16)/(n(n-1))\) on its orthogonal complement. Therefore

\[
\|\operatorname{Cov}_{\nu_n}(X)\|_{\rm op}\le2\quad(n\ge8\text{ even}).
\tag{3.1}
\]

At \(n=6\) the norm is \(8/3\), not two. The same norm bound two holds
for \(\mu _n\) when \(n\ge 8\), since its covariance is
\((1-e^{-n})\operatorname{Cov}_\nu +e^{-n}I\). For \(M=S/n\),

\[
\operatorname{Var}_{\mu_n}(M)=(1-e^{-n})\frac{16}{n^2}+e^{-n}\frac1n
\sim\frac{16}{n^2}.
\tag{3.2}
\]

Moreover \(\mu _n(|M|=4/n)\ge 1-e^{-n}\). These distributions have
stronger magnetization concentration than the uniform cube, bounded
covariance norm, and no constant-variance distributed linear witness,
but still fail full polynomial AA.

Bounded covariance does imply a degree-one theorem: if
\(\lVert \operatorname{Cov}(X)\rVert_{\mathrm{op}}\le C\) and
\(f=a_{0}+\sum_i a_iX_i\) is globally bounded by one, then
\(\sum_i|a_i|\le 1\) and \(\operatorname{Var}(f)\le C \max _i|a_i|\).
Thus \(\max _i I_i(f)\ge \operatorname{Var}(f)^{2}/C^{2}\). Theorem T1
shows why this does not automatically extend to higher fixed degrees.
